\documentclass[10pt,a4paper]{amsart}
\usepackage[margin=19mm]{geometry}
\usepackage{amsmath,amssymb,mathtools}
\usepackage[scr=rsfs,cal=euler]{mathalfa}
\usepackage{xfrac}
\usepackage[unicode,hidelinks,bookmarks=false]{hyperref}
\newcommand{\ie}{i.e.\ }
\newcommand{\cf}{cf.\ }
\newcommand{\resp}{respectively\ }
\newcommand{\Subsection}{\subsection}
\providecommand{\nobreakdash}{\nobreak\hbox{-}}
\setlength{\emergencystretch}{2em}
\multlinegap0pt
\usepackage{tikz}
\usetikzlibrary{cd}
\usepackage{amssymb}
\usepackage{enumerate}
\usepackage{cancel}
\usepackage{xcolor}
\usepackage{tikz}
\newtheorem{prop}{Proposition}
\newtheorem{lem}{Lemma}
\newtheorem{thm}{Theorem}
\DeclareMathOperator{\Aut}{Aut}
\DeclareMathOperator{\Image}{Im}
\newcommand{\ra}{\rightarrow}
\newcommand{\tC}{{\widetilde C}}
\title{Finiteness and injectivity of Prym maps for cyclic coverings}
\author{Pawe{\l} Bor\'owka, Juan Carlos Naranjo, Angela Ortega, Anatoli Shatsila}
\date{Source: arXiv:2510.01330v3 (CC BY 4.0)}
\begin{document}
\maketitle

\noindent\emph{Source and licence.} Finiteness and injectivity of Prym maps for cyclic coverings. Version arXiv:2510.01330v3 (CC BY 4.0), distributed by arXiv under the Creative Commons Attribution 4.0 International licence, http://creativecommons.org/licenses/by/4.0/, as recorded in the arXiv licence metadata for this exact version and shown on the abstract page https://arxiv.org/abs/2510.01330v3. Full licence notice: \texttt{licenses/arxiv-2510.01330-CC-BY-4.0.txt}.\par\smallskip

\noindent\textit{Editorial scope.} Counted unit: the article's thm mainodd (source0.tex lines 503--505). The article's complete original proof of it is delivered with it (source0.tex lines 507--524, one \texttt{proof} environment of 18 lines). No mathematics is written, repaired or re-derived: the statement and the proof are the article's own lines, in the article's own notation, and no retained line is substituted or re-flowed.  Retained uncounted as the source-local context the proof needs: lines 207--212; lines 270--272; lines 282--285; lines 433--443; lines 452--454; lines 484--490.  Nothing was omitted from any retained span, so the package carries no omission record.  No retained line contains a citation, so the wrapper carries no bibliography and no bibliography entry.  The retained lines contain the article's own diagram code, which the wrapper typesets with the standard tikz packages; no image file is delivered and the package declares no asset dependency.  Editorial formatting only, in addition: an AMS \texttt{amsart} preamble that restates the article's own declarations and macros used by the retained lines, namely the environments \texttt{prop}, \texttt{lem}, \texttt{thm} on separate counters, together with the standard packages those lines need.  The retained environments print in this document as Proposition 1, Lemma 1, Theorem 1 in order of appearance, whereas the article prints the same items with the numbering of its own full document.  The complete native source of the article as distributed in the arXiv e-print for this exact version is bundled unchanged as \texttt{sources/186198/source0.tex}.
\begin{prop}
\label{group}
    Let $d \geq 3$ be an integer and $(P, \Theta_P)$ be an element of $\Image\mathcal{P_H}_g[d]$. 
    Then the subgroup of automorphisms $$G := \{ \phi \in \Aut(P,\Theta_P) \: | \: \phi(x) = x \:\: \forall x \in K(\Theta_P)\}.$$
    is isomorphic to $\langle \sigma, -j \rangle \simeq {D}_{d}$.
\end{prop}

\begin{lem}\label{embeddingE}
Let $E_0 := \tC/\langle\sigma^p,j\rangle$ and $\pi: \tC \to E_0$ be the quotient map. Then $\pi^*:JE_0\to J\tC$ is an embedding. In particular, $\pi:\tC\to E_0$ does not factorise via $\tC\to C''\to E_0$ where $C''\to E_0$ is cyclic \'etale. 
\end{lem}

\begin{lem}\label{key_lemma}
    Let $h:C'\to C$ be a covering that does not factorise via an \'etale covering $C'\to C''\to C$.
        Then $h^*: JC\to JC'$ is an embedding and any other embedding $i:JC\to JC'$ with $\Image i=\Image h^*$ is (up to an isomorphism of $JC$) given by $h^*$.
\end{lem}

\begin{equation}
\label{small-tower}
% https://tikzcd.yichuanshen.de/#N4Igdg9gJgpgziAXAbVABwnAlgFyxMJZABgBpiBdUkANwEMAbAVxiRAB12cBhEAX1LpMufIRQAmUuKq1GLNtwD6wTtgDmAWzqKAzH36CQGbHgJEyOmfWatEIABIGhJ0UR2kAjFbm2QS4k5GwqZiyAAspJbU1vJ2AKKKAQLOImYSpGHeNmycWjgAFgBGhcAACnwAeh78MjBQavBEoABmAE4QGkhkIDgQSB7RPmxQiNXJIG0d-dS9SJKy2XY6o4GTnYjzs4gArIOLIOIr42tIET19O3uxIMtjhieI7uenV74jdy3t691bTzG+hw+Ey+XRmF3m-zYOneqxBozBLwW10BNT4QA
\begin{tikzcd}
\tC \arrow[rrdd, "q:1"] \arrow[rrrd, "2:1"] \arrow[ddd, "d:1"] &  &                                                    &                        &                        \\
                                                                &  &                                                    & C_0 \arrow[rdd, "q:1"] &                        \\
                                                                &  & C_{\sigma^p} \arrow[lld, "p:1"] \arrow[rrd, "2:1"] &                        &                        \\
H \arrow[rrd, "2:1"]                                            &  &                                                    &                        & E_0 \arrow[lld, "p:1"] \\
                                                                &  & \mathbb{P}^1                                       &                        &                       
\end{tikzcd}
\end{equation}

\begin{lem}\label{elliptic2}
   The covering $C_0\to E_0$ of degree $q$ is never Galois. Its Galois closure is $\tC\to E_0$. 
\end{lem}

\begin{lem}\label{lemunique}
    %Let $Q\in\Image\mathcal{P_H}_g[p]$ be general.
    Let $p$ be an odd prime number.
    For a general $[h:C\to H]\in \mathcal{RH}_g[p]$, %\mathcal{P_H}_g[p]^{-1}(Q)$ fix the curve $C$ and define $X_C:=\{[h':C\to H']\in\mathcal{P_H}_g[p]^{-1}(Q)\}$. Then $X_C=\{[h:C\to H]\}$ is a singleton and 
    there is no other $h':C'\to H'$ with $C'\simeq C$ and $P(h)\simeq P(h')$. 
    In particular, the embedding $P(h)\subseteq JC$ is unique (up to an automorphism of $P(h)$).
\end{lem}

\begin{thm}\label{mainodd}
Let $d$ be an odd non-prime number. Then, the Prym map  $$\mathcal{P_H}_g[d]: \mathcal{RH}_g[d]\to\mathcal{A}_{(g-1)(d-1)}^{(1,1,\ldots, 1,d,\dots,d)}$$ is generically injective.
\end{thm} 

\begin{proof}
	Let $(P,\Theta_P)$ be an element of the image of $\mathcal{P}_{d}$. By Proposition \ref{group} we recover $\langle \sigma, -j  \rangle \subset \Aut(P,\Theta_P)$ as a subgroup isomorphic to $D_{d}$. 
    Note that $\mathbb{Z}_{d} \simeq \langle \sigma \rangle \subsetneq \langle \sigma, -j \rangle \simeq D_{d}$ is the unique subgroup of $D_{d}$ isomorphic to $\mathbb{Z}_{d}$, hence we can choose $\sigma \in \Aut(P,\Theta_P)$ to be one of the generators of $\mathbb{Z}_{d}$. 
    The elements of $D_{d} \setminus \mathbb{Z}_{d}$ are involutions of the form $-j\sigma^s$, where $s = 0,\ldots, d-1$, so let us fix one them which we denote by $-j$. Let $A := \Image(1+j)$ and  $B := \Image(1+j)(1+\sigma^p+\ldots+\sigma^{p(q-1)})$. 
    
Let $f:\tC\to H$ be an element in the preimage $\mathcal{P_H}_g[d]^{-1}(P)$. Then, by Proposition \ref{group} we have $j\in \Aut(\tC)$ is an involution with fixed points so $A=\Image (1+j)=JC_0$ is the Jacobian of the quotient curve $C_0=\tC/j$ embedded in $P$. Similarly, $B=JE_0$ by Lemma \ref{embeddingE}.


   Since $JE_0 \subsetneq JC_0$, we can use Lemma \ref{key_lemma} to get a unique covering $C_0 \ra E_0$ which recovers the covering map of degree $q$ given in Diagram \eqref{small-tower}.

    By Lemma \ref{elliptic2} the Galois closure of $C_0 \to E_0$ is precisely the covering $\tC \to E_0$ and it comes  with an automorphism $\sigma^p$ of order $q$ and an involution $j$ on $\tC$. Dividing by $\sigma^p$, we get a curve $C_{\sigma^p}$. 
    
       Denote by $g:\tC\to C_{\sigma^p}$ and $h:C_{\sigma^p}\to H$ the quotient maps, where $h$ is a quotient of $C_{\sigma^p}$ by the push-down of $\sigma\in\Aut(\tC)$. Note that $J\tC=f^*(JH)\boxplus 
       P(f)$ with $P\cong P(f)$ and $g^*JC_{\sigma^p}=g^*(h^*(JH))\boxplus g^*(P(h))$. However, in the proof of Lemma \ref{embeddingE} we have shown $g^*_{|P(h)}$ is an embedding, so $g^*JC_{\sigma^p}\cap P(f)=g^*P(h)\cong P(h)$. 

    Now, to sum up, from the Prym variety $P$ we have recovered $C_{\sigma^p}$ and $JE_0\times JE_0$ which is isomorphic to the  Prym variety of some covering $h$ of degree $p$. Hence,  assuming $P$ is general, according to Lemma \ref{lemunique}, $h$ is uniquely defined by the construction and hence $f=h\circ g$ is unique.
    Therefore, the Prym map is generically injective. 
    \end{proof}

\end{document}
